% Formatação extraída do template COTB 2027 fornecido.
% Compatibilidade: language=english por último mantém inglês como padrão
% do babel, sem a chave aninhada main=english que falha em xkeyval antigo.
\documentclass[sigconf, nonacm, screen, language=portuguese, language=spanish, language=english]{acmart}             % VERSAO FINAL

\AtBeginDocument{%
  \providecommand\BibTeX{{%
    \normalfont B\kern-0.5em{\scshape i\kern-0.25em b}\kern-0.8em\TeX}}}

%% NAO carregue o pacote babel manualmente: a classe acmart ja o
%% carrega automaticamente com os idiomas indicados nas opcoes
%% "language=..." do \documentclass acima (evita conflito de dupla
%% carga do babel).
\usepackage{amsmath}
\usepackage{amsfonts}
\let\Bbbk\relax
\usepackage{amssymb}
\usepackage{booktabs}
\usepackage{multirow}
\usepackage{graphicx}

%% Remove o bloco de identificacao ACM (nao se aplica ao evento)
\settopmatter{printccs=false, printacmref=false}
\renewcommand\footnotetextcopyrightpermission[1]{}

%% =====================================================================
%% CABECALHO PERSONALIZADO (dados do evento a esquerda, autores a
%% direita)
%% ---------------------------------------------------------------------
%% - O lado ESQUERDO do cabecalho traz o nome/edicao do evento, local
%%   e datas. Edite a macro \cotbEventInfo abaixo se necessario.
%% - O lado DIREITO do cabecalho traz os autores do trabalho:
%%     * ate 3 (tres) autores: lista todos os nomes (gerado
%%       automaticamente a partir dos comandos \author usados mais
%%       abaixo, no mesmo formato usado pela classe acmart);
%%     * mais de 3 (tres) autores: exibe apenas o PRIMEIRO autor
%%       seguido de "et al." (tambem gerado automaticamente).
%%   Nao e necessario editar nada manualmente: o comportamento eh
%%   determinado pela quantidade de comandos \author preenchidos
%%   na secao de autores, logo abaixo.
%% =====================================================================
\makeatletter
\newcounter{cotbNumAuthors}
\let\cotb@originalauthor\author
\renewcommand{\author}[1]{%
  \stepcounter{cotbNumAuthors}%
  \ifnum\value{cotbNumAuthors}=1
    \gdef\cotbFirstAuthorName{#1}%
  \fi
  \cotb@originalauthor{#1}%
}
\makeatother

%% Dados do evento exibidos do lado esquerdo do cabecalho (mantenha
%% o texto curto, em ate duas linhas, para nao colidir com o lado
%% direito):
\newcommand{\cotbEventInfo}{XVIII Computer on the Beach (COTB 2027)\\
  Florianópolis/SC --- 13 a 17 de abril de 2027}

%% Autores exibidos do lado direito do cabecalho (logica automatica
%% de "et al." para mais de 3 autores):
\newcommand{\cotbHeaderAuthors}{%
  \ifnum\value{cotbNumAuthors}>3
    \cotbFirstAuthorName\ et al.%
  \else
    \shortauthors
  \fi
}

%% Espaco reservado para o cabecalho de duas linhas nao colidir com
%% o corpo do texto:
\setlength{\headheight}{24pt}

\AtBeginDocument{%
  \fancypagestyle{firstpagestyle}{%
    \fancyhf{}%
    \renewcommand{\headrulewidth}{0.4pt}%
    \renewcommand{\footrulewidth}{0pt}%
    \fancyhead[L]{\begin{minipage}[b]{0.62\textwidth}\raggedright
      \scriptsize\cotbEventInfo\end{minipage}}%
    \fancyhead[R]{\begin{minipage}[b]{0.34\textwidth}\raggedleft
      \scriptsize\cotbHeaderAuthors\end{minipage}}%
    \fancyfoot[C]{\if@ACM@printfolios\thepage\fi}%
  }%
  \fancypagestyle{standardpagestyle}{%
    \fancyhf{}%
    \renewcommand{\headrulewidth}{0.4pt}%
    \renewcommand{\footrulewidth}{0pt}%
    \fancyhead[L]{\begin{minipage}[b]{0.62\textwidth}\raggedright
      \scriptsize\cotbEventInfo\end{minipage}}%
    \fancyhead[R]{\begin{minipage}[b]{0.34\textwidth}\raggedleft
      \scriptsize\cotbHeaderAuthors\end{minipage}}%
    \fancyfoot[C]{\if@ACM@printfolios\thepage\fi}%
  }%
  %% IMPORTANTE: a classe acmart ja executa "\pagestyle{standardpagestyle}"
  %% durante o carregamento (antes desta redefinicao ter efeito). Por
  %% isso, e necessario chamar \pagestyle novamente aqui, para que o
  %% cabecalho das paginas internas passe a usar a versao customizada
  %% definida acima (evento a esquerda, autores a direita).
  \pagestyle{standardpagestyle}%
}
